\subsection*{Notation and convention}
In this article, we only work with smooth manifolds. Given a vector bundle $A$ over a smooth manifold $M$, a {\em distribution} is an assignment $m\in M\mapsto R_m\subseteq A_m$, where $R_m$ is a vector subspace of the fiber $A_m$; when any $v\in R_m$ can be extended to a local section of $A$ taking values in $R$, we say that the distribution is {\em smooth}. Each distribution has associated a $\mathbb{N}$-valued function given by the assignment $\rk\colon m\in M\mapsto\dim R_m$, which is called {\em rank}. A distribution is said to be {\em regular} if it is a vector bundle.

Given a distribution $E\subset A_\C$, where $A_{\C}$ is the complexification of a real vector bundle $A\to M$, we denote the space of real elements of $E$ by $\re E:=E\cap A$ and call them the {\em real part of $E$}. Given a smooth map $\varphi\colon M\to N$, we denote the complexification of $T\varphi$ by $T_{\C}\varphi$.

Throughout this article, a real anchored vector bundle is referred to by RAVB and a real Lie algebroid by RLA. In Appendix~\ref{preliminaries0}, we recall some properties of RAVBs and RLAs.

Let $N$ be a submanifold of $M$, the normal bundle of $N$ is denoted by
\[
\nu(M,N)=TM|_{N}/TN,
\]
when the context is clear we denote it by $\nu_N$. In Appendix \ref{preliminarieslocstr}, we recall some properties of normal vector bundles.

\section{Complex anchored vector bundles}\label{CAVB}
\subsection{Complex anchored vector bundles}
A complex vector bundle $E\to M$ is equipped with a canonical bundle map $j_E\in \Diff(E)$ defined by $j_E(e)={\rm i}\cdot e$, the multiplication by the imaginary number ${\rm i}$. The {\em group of $($complex$)$ automorphisms of $E$} is
 \[
 \Aut(E)=\{(\Phi,\varphi)\in \Aut(E_{\R})\mid \Phi\circ j_E=j_E \circ \Phi\ \text{and}\ \Phi\ \text{covers}\ \varphi\in \Diff(M)\},
 \]
 where $E_{\R}$ denotes the realification of $E$, and $\Aut(E_{\R})$ is the automorphism group of $E_{\R}$ as a~real vector bundle.
The map $j_E$ acts on $\Aut(E)$ in the natural way. Consider a 1-parameter subgroup $\{\Phi_t\}_{t\in \R}$ of $\Aut(E)$. Since $\Phi_t\in \Aut(E_{\R})$ and $\Phi_t\circ j_E=j_E\circ \Phi_t$, the vector field \smash{$\hat{X}=\frac{{\rm d}}{{\rm d}t}\Phi_t|_{t=0}\in \mathfrak{aut}(E_{\R})$} satisfies \smash{$(j_E)_* \hat{X}=\hat{X}$}. So, the {\em Lie algebra of $($complex$)$ infinitesimal automorphisms of $E$} is given by
 \[\mathfrak{aut}(E)=\bigl\{\hat{X}\in \mathfrak{aut}(E_{\R})\mid (j_E)_* \hat{X}=\hat{X}\bigr\}\subseteq \mathfrak{aut}(E_{\R})\subseteq \mathfrak{X}(E).\]
 Another characterization of $\aut(E)$ is the following.

 \begin{Lemma} The Lie algebra $\mathfrak{aut}(E)$ is given by the pairs $(L,X)$ such that $X\in \mathfrak{X}(M)$ and $L\colon \Gamma (E)\to \Gamma (E)$ is a $\C$-linear derivation with respect to $X$.
 \end{Lemma}
\begin{proof}
Since $\widetilde{X}\in \aut(E)\subseteq \aut (E_{\R})$, it has associated a derivation given by the pair $(L, X)$, where $L\colon \Gamma(E_{\R})\to \Gamma(E_{\R})$ is $\R$-linear and \smash{$X=\pr_* \widetilde{X}\in \X(M)$}, defined as \smash{$L(\sigma)^{\uparrow}=[\sigma^{\uparrow}, \widetilde{X}]$}, ${\forall \sigma\in \Gamma(E)}$, where $\uparrow$ denotes the vertical lift, see Appendix \ref{preliminaries0}. Note that $L$ is $\C$-linear if and only if $L({\rm i}\sigma)={\rm i}L(\sigma)$ and thus
\begin{align*}
L({\rm i}\sigma)^{\uparrow}&{}=\bigl[({\rm i}\sigma)^{\uparrow}, \widetilde{X}\bigr]=\bigl[(j_{E})_* \sigma^{\uparrow}, \widetilde{X}\bigr]=\bigl[(j_{E})_* \sigma^{\uparrow},(j_{E})_* \widetilde{X}\bigr]\\
&{}=(j_E)_* \bigl[\sigma^{\uparrow}, \widetilde{X}\bigr]=(j_{E})_* \bigl(D(\sigma)^{\uparrow}\bigr)=({\rm i}L(\sigma))^{\uparrow}.
\end{align*}
And by injectivity of the ``$\uparrow$''-map, we have that $L({\rm i}\sigma)={\rm i}L(\sigma)$ and the lemma holds.
\end{proof}

\begin{Definition}
 A {\em complex anchored vector bundle (CAVB)} is a pair $(A,\rho)$, where $A$ is a complex vector bundle over $M$, and $\rho\colon A\to T_{\C}M$ is a $\C$-linear bundle map. Let $(A,\rho_A)$ and $(B,\rho_B)$ be two CAVBs over $M$ and $N$, respectively. A {\em morphism of complex anchored vector bundles} consists of a bundle map $\Phi\colon A\to B$ and a map $\varphi\colon M\to N$ such that the following diagram commutes:
\[\xymatrix{A\ar[d]^{\rho_A}\ar[r]^{\Phi} & B\ar[d]^{\rho_B}\\ T_{\C}M\ar[r]^{T_{\C}\varphi} & T_{\C}N.
}\]

\end{Definition}
Given a CAVB $(A, \rho)$, the anchor map decomposes $\rho=\rho_1+{\rm i}\rho_2$. So we obtain two real anchored vector bundles $(A_{\R},\rho_1)$ and $(A_{\R}, \rho_2)$. By the $\C$-linearity of $\rho$, we have that
\[\rho_1 \circ j_A=-\rho_2.\]
\begin{Proposition}
The map $j_A$ is an isomorphism of RAVBs between $(A, \rho_1)$ and $(A,-\rho_2)$.
\end{Proposition}
 The {\em group of automorphisms of a CAVB} $(A,\rho)$ is given by
 \[\Aut(A,\rho)=\{(\Phi, \varphi)\in \Aut(A)\mid \rho\circ\Phi=T_{\C}\varphi\circ \rho\}.\]
 Note that $\rho\circ\Phi=T_{\C}\varphi\circ \rho$ if and only if $\rho_1\circ \Phi=T\varphi\circ \rho_1$ and $\rho_2\circ \Phi=T\varphi\circ \rho_2$. So
 \[\Aut(A,\rho)=\Aut(A_{\R},\rho_1)\cap \Aut(A_{\R},\rho_2)\cap \Aut (A).\]
 The {\em $($complex$)$ infinitesimal automorphisms of $(A,\rho)$} are denoted by $\aut(A,\rho)$.
\begin{Definition}\label{cx_tangent_extension}
   The {\em complex tangent extension} of a vector field $X\in \X(M)$, is the vector field $X_{T_{\C}}\in \mathfrak{X}(T_{\C}M)$ defined by the one-parameter subgroup \smash{$T_{\C}\varphi_t ^{X}$}, where \smash{$\varphi_t ^{X}$} is the local flow of $X$.
\end{Definition}

For a CAVB $(A,\rho)$, we have that
\begin{align*}
\mathfrak{aut}(A,\rho)&{}=\aut(A_{\R},\rho_1)\cap \aut(A_{\R},\rho_2)\cap \aut(A)\\
&{}=\bigl\{\widetilde{X}\in \X(A)\mid (j_E)_* \hat{X}=\hat{X},\, \widetilde{X}\sim_{\rho} X_{T_{\C}},\, \text{where}\ \widetilde{X}\sim_{\pr} X\bigr\}.
\end{align*}



In terms of the differential operators: a pair $(L,X)$, where $L\colon \Gamma(A)\to \Gamma(A)$ is a $\C$-linear differential operator and $X\in \mathfrak{X}(M)$, belongs to $\mathfrak{aut}(A,\rho)$ if and only if $\rho(L(\tau))=[X,\rho(\tau)]$
or equivalently
\begin{equation*}%\label{prop_aut_AVB}
\rho_i(L(\tau))=[X, \rho_i(\tau)]\qquad\text{for $i=1,2$}.
\end{equation*}
\begin{Definition}
A real or complex anchored vector bundle $(A,\rho)$ is {\em involutive} if $\rho(\Gamma(A))$ is a~Lie subalgebra of $\Gamma(TM)$ or $\Gamma(T_{\C}M)$, respectively.
\end{Definition}

Consider the decomposition of the anchor map $\rho$ in its real and imaginary parts ${\rho=\rho_1+{\rm i}\rho_2}$. Note that the involutivity of $(A,\rho)$ is not directly related to the involutivity of $(A_{\R}, \rho_1)$ and $(A_{\R}, \rho_2)$. Even if both $(A_{\R}, \rho_1)$ and $(A_{\R}, \rho_2)$ are involutive, we cannot ensure the involutivity of~$(A,\rho)$ and vice versa. To solve this issue, we shall introduce the distribution of real elements in the next subsection.

\subsection{Distributions associated to CAVBs}
A CAVB $(A,\rho)$ has associated two real distributions~in~$TM$:
 \[\Delta=\re \bigl(\rho(A)\cap \overline{\rho(A)}\bigr)\qquad\text{and} \qquad D=\re \bigl(\rho(A) + \overline{\rho(A)}\bigr).\]
Observe that $D$ is always a smooth distribution, while $\Delta$ is not necessarily smooth. However, we have the following.

 \begin{Lemma}\label{imageanchor}
 If $D$ is a regular distribution, then $\Delta$ is a smooth distribution.
 \end{Lemma}
\begin{proof}
 Consider the map $\rho_2\colon A\to TM$, which is a real bundle map. Since the image of $\rho_2$ is~$D$, its kernel $\ker \rho_2$ is a vector bundle. Thus, $\Delta=\rho(\ker \rho_2)$ is smooth.
\end{proof}


\begin{Definition}
Given a CAVB $(A,\rho)$ define the {\em distribution of real elements of $(A,\rho)$} as
\[A^{\real} =\{e\in A\mid \rho_2(e)=0\},\]
which we called {\em distribution of real elements of $(A,\rho)$}. In general, $A^{\real}$ is a distribution in $A$ (not necessarily smooth). When $A^{\real}$ is regular, the pair $(A^{\real},\rho^{\real})$ is a RAVB, where $\rho^{\real}=\rho_1|_{A^{\real}}$. The $\mathbb{N}$-valued function % \text{real-rank A} -> \operatorname{\text{real-rank}} A
\[\operatorname{\text{real-rank}} A\colon \ m\in M\mapsto \dim_{\R} A^{\real}|_m\]
is called {\em real rank of~$A$}.
\end{Definition}
Constant real rank is equivalent to $A^{\real}$ being a vector bundle and by the proof of Lemma \ref{imageanchor}, it is equivalent to the regularity of $D$ as a distribution. Indeed, we can see that
\begin{equation}\label{sum_identity}
  {\operatorname{\text{real-rank}} A}+\rk_{\R} D=\rk_{\R} A_{\R}.
\end{equation}
\begin{Lemma}\label{restriction_inf_aut}
   Assume that $(A,\rho)$ has constant real rank. If $(L,X)\in \aut(A,\rho)$, then $L(\Gamma(A^{\real}))\subseteq\Gamma(A^{\real})$ and $(L|_{A^{\real}}, X)\in \aut(A^{\real},\rho^{\real})$.
\end{Lemma}
\begin{proof}
  Straightforward.
\end{proof}

 \begin{Lemma}\label{cavb_inv}
 If $(A,\rho)$ is an involutive CAVB with constant real rank, then $(A^{\real}, \rho^{\real})$ is involutive.
 \end{Lemma}
\begin{proof}
Let $\alpha, \beta \in \Gamma (A^{\real})$, equivalently, $\rho (\alpha)$ and $\rho (\beta)$ are real, and so $[\rho (\alpha), \rho (\beta)]$. By involutivity of $A$, there exists $\gamma \in \Gamma (A)$ such that $[\rho (\alpha), \rho (\beta)]=\rho (\gamma)$. So $\rho (\gamma)$ is real and consequently~${\gamma\in \Gamma(A^{\real})}$.
\end{proof}


The image of the anchor map of an involutive RAVB is integrable, see for example \cite{bursztyn2019splitting}. In~the~case of a CAVB $(A,\rho)$, it follows from the proof of Lemma \ref{imageanchor} and by Lemma \ref{cavb_inv} that the real part of $\rho(A)$ is integrable.


The pullback of complex anchored vector bundle $(A, \rho) $ over $M$ along the map $\varphi\colon N\to M$ is defined in the same way as the real case, as the fibered product:
\[
\xymatrix{\varphi^! A\ar[d]\ar[r] & A \ar[d]^{\rho}\\
T_{\C}N\ar[r]^{T_{\C}\varphi} & T_{\C}M.}
\]

Note that
\begin{align*}
\varphi^! A&{}=\bigl\{(a, X_1+{\rm i}X_2) \mid  (a,X_1)\in \varphi^! A_1\ \text{and}\ (a,X_2)\in \varphi^! A_2\bigr\}\\
&{}=\varphi^! A_1 \times_{\pr_A} {\rm i}\varphi^! A_2\subseteq A\times T_{\C}N,
\end{align*}
where $A_1=(A_{\R}, \rho_1)$ and $A_2=(A_{\R}, \rho_2)$, and $\pr_A\colon A\times T_{\C}N\to A$ is the obvious projection.
As~a~consequence, we can see that
\begin{enumerate}\itemsep=0pt
\item $\varphi^!T_{\C}M=T_{\C}N$.
\item If $N=M\times Q$, then $\pr_M ^! A=A\times T_{\C}Q$.
\item If \smash{$N\xhookrightarrow{\iota} M$} is a submanifold transversal to $\rho$, then $N$ is also transversal to $\rho_1$ and $\rho_2$. Moreover,
\[
\bigl(\iota^! A\bigr)_{\R}=\rho^{-1}(T_{\C}N)_{\R}=\{a\in A_{\R}\mid \rho(a)=\rho_1(a)+{\rm i}\rho_2(a)\in T_{\C}N\}=\iota^{!}A_1\cap \iota^{!}A_2.
\]
\item If $\rho$ is injective, then $\varphi^! A=(T_{\C}\varphi)^{-1}(A)$.
\end{enumerate}

 The following lemma is straightforward.

\begin{Lemma}\label{repartbackward}
Let $(A,\rho)$ be a CAVB and a map $\varphi\colon N\to M$. Then,
\begin{align*} % RS: im -> \im
&\bigl(\varphi^! A\bigr)^{\real}=\varphi^! (A^{\real})\times_{\pr_A}{\rm i}(A\times \ker T\varphi),\\
&\Delta_{\varphi! A}=\im(T\varphi)\cap \Delta_{A}.
\end{align*}
\end{Lemma}

A submanifold \smash{$N\xhookrightarrow{\iota} M$} is called {\em transversal to the CAVB} if
\[\rho(A)+T_{\C}N=T_{\C}M\]
(or the stronger condition $\Delta+TN=TM$). By transversality, \smash{$\iota^! A$} is smooth and so is a CAVB. Note that transversality with respect to $\rho$ implies transversality with respect to $A_1$ and $A_2$.

A direct consequence of the Lemma \ref{repartbackward} is that \smash{$\bigl(\iota^! A\bigr)^{\real}=\iota^! (A^{\real})$} for submanifolds \smash{$N\xhookrightarrow{\iota} M$} (not necessarily transversal to the anchor map).
The following lemma appeared originally in~\cite{bursztyn2019splitting}, and the proof is similar in the complex case.
\begin{Lemma}
   Suppose $(A,\rho)$ is an involutive CAVB over $M$ and $\varphi\colon N\to M$ a map transversal to $\rho$. Then $\varphi^! A$ is an involutive CAVB.
\end{Lemma}

 \section{Complex Lie algebroids}\label{CLA}
\begin{Definition}%\label{def_cla}
A {\em complex Lie algebroid $($CLA$)$} over a manifold $M$ is a complex anchored vector bundle $(A, \rho)$ over $M$ together with a Lie bracket on $\Gamma(A)$ satisfying the Jacobi identity
\[[\cdot,\cdot]\colon \ \Gamma(A)\times \Gamma(A)\to\Gamma(A)\]
 and with a $\C$-bundle map $\rho\colon A\to T_{\C}M$, that we call the anchor map, preserving brackets and satisfying the Leibniz identity
 \[[\alpha,f \beta]=f[\alpha,\beta]+\rho(\alpha)(f)\beta\]
 for all $\alpha,\beta\in \Gamma(A)$ and for any function $f\in C^{\infty}(M,\C)$. A {\em complex skew-algebroid} is when the Jacobi identity is not necessarily satisfied, and an {\em almost complex Lie algebroid $($almost CLA$)$} is a complex skew-algebroid where the anchor preserves the brackets.
\end{Definition}

\begin{Remark}
Observe that CLAs are defined on the category of smooth real manifolds and not on the category of holomorphic manifolds.
\end{Remark}
Morphisms between CLAs are defined in the exact same way as RLAs, and so are their automorphisms.
